\pdfminorversion=7%
\documentclass[aspectratio=169,mathserif,notheorems]{beamer}%
%
\xdef\bookbaseDir{../../bookbase}%
\xdef\sharedDir{../../shared}%
\RequirePackage{\bookbaseDir/styles/slides}%
\RequirePackage{\sharedDir/styles/styles}%
\toggleToGerman%
%
\subtitle{35.~PgModeler Installieren}%
%
\begin{document}%
%
\startPresentation%
%
\section{Einleitung}%
%
\begin{frame}%
\frametitle{PgModeler}%
\bigskip\strut\bigskip%
\begin{itemize}%
\item Um konzeptuelle Modelle zu erstellen, hatten wir \yEd\ verwendet.%
\item<2-> Konzeptuelle Modelle sind an keine Technologie gebunden, daher war ein allgemeiner Graphen-Editor wie \yEd\ genau richtig.%
\item<3-> Logische Modelle sind oftmals an eine konkrete Datenbanktechnologie wie \postgresql\ gebunden.%
\item<4-> Wir können sie in \glsShort{SQL} definieren {\dots} aber es gibt auch Werkzeuge, mit denen wir sie über eine \glsShort{GUI} erstellen können.%
\item<5-> Wir verwenden den \pgmodeler\cite{AES2006PPDM}, mit dem logischen Modelle für \postgresql\ über eine \glsShort{GUI} entworfen werden können.%
\end{itemize}%
\locateGraphic[%
The \href{https://pgmodeler.io}{\pgmodeler\ logo} is under the copyright of Raphael Araújo~e~Silva.%
]{}{width=0.2\paperwidth}{../05_software_und_literatur/graphics/pgmodelerLogo}{0.75}{0.02}%
\end{frame}%
%
\section{PgModeler unter Ubuntu Linux Installieren}%
%
\begin{frame}[t]%
\frametitle{PgModeler unter Ubuntu Linux Installieren}%
%
\begin{itemize}%
%
\only<-1>{\item Installieren wir nun den \pgmodeler\ unter \ubuntu\ \linux.}%
%
\only<2-3>{%
\item<2-> Öffnen Sie ein \bash\ \glslink{terminal}{Terminal} via with \ubuntuTerminal.%
\item<3-> Schreiben Sie \bashil{sudo apt-get install pgmodeler} und drücken~\keys{\return}.%
}%
%
\only<4>{%
\item<4-> Wenn Sie nach dem \pgls{sudo}-Passwort gefragt werden, tippen Sie es ein und drücken~\keys{\return}.%
}%
%
\only<-5>{%
\item<5-> Das System sagt uns, dass Pakete installiert werden müssen und fragt uns, ob das OK ist.%
}\only<-6>{%
\item<6-> Das ist OK. Wir drücken~\keys{Y} und dann~\keys{\return}.%
}%
%
\only<7>{%
\item<7-> Der \pgmodeler\ wird installiert.%
}%
%
\only<-8>{%
\item<8-> Wir führen den \pgmodeler\ aus, in dem wir \bashil{pgmodeler} in das \bash\ \glslink{terminal}{Terminal} eintippen und \keys{\return}~drücken.%
}%
%
\only<-10>{%
\item<9-> Das Programm zeigt vielleicht ein Fehler-Fenster an.%
\item<10-> Wenn das passiert, können wir das einfach ignorieren und auf \menu{OK} klicken.%
}%
%
\only<11>{%
\item<11-> Das \pgmodeler-Fenster öffnet sich.%
}%
%
\only<-13>{%
\item<12-> Ich verwende \pgmodeler\ im \inQuotes{Light Mode}.%
\item<13-> Um den zu aktivieren, klicken wir auf \menu{\pgmodelerMainMenu>Edit>Settings}.%
}%
%
\only<14>{%
\item<14-> Im Menü \menu{Appearance} ändern wir das \inQuotes{Theme} zu \menu{Light}{\dots}%
}%
%
\only<15>{%
\item Wir klicken auf \menu{Apply}.%
}%
%
\only<16>{%
\item<16-> Im Register \menu{Relationships} stellen wir sicher, dass der \inQuotes{connection mode} auf \inQuotes{Crow's foot notation} gesetzt ist.%
}%
%
\item<17-> Damit haben wir den \pgmodeler\ installiert.%
%
\end{itemize}%
%
\locateGraphic{2-3}{width=0.64\paperwidth}{graphics/installingPgModelerUbuntu/installingPgmodelerUbuntu01aptGetInstall}{0.18}{0.33}%
\locateGraphic{4}{width=0.64\paperwidth}{graphics/installingPgModelerUbuntu/installingPgmodelerUbuntu02pw}{0.18}{0.33}%
\locateGraphic{5}{width=0.4\paperwidth}{graphics/installingPgModelerUbuntu/installingPgmodelerUbuntu03yn}{0.3}{0.21}%
\locateGraphic{6}{width=0.4\paperwidth}{graphics/installingPgModelerUbuntu/installingPgmodelerUbuntu04yny}{0.3}{0.21}%
\locateGraphic{7}{width=0.64\paperwidth}{graphics/installingPgModelerUbuntu/installingPgmodelerUbuntu05installed}{0.18}{0.33}%
\locateGraphic{8}{width=0.64\paperwidth}{graphics/installingPgModelerUbuntu/installingPgmodelerUbuntu06run}{0.18}{0.33}%
\locateGraphicTB{9-10}{width=0.64\paperwidth}{graphics/installingPgModelerUbuntu/installingPgmodelerUbuntu07error}{0.18}{0.3}%
\locateGraphicTB{11}{width=0.64\paperwidth}{graphics/installingPgModelerUbuntu/installingPgmodelerUbuntu08open}{0.18}{0.33}%
\locateGraphicTB{12-13}{width=0.55\paperwidth}{graphics/installingPgModelerUbuntu/installingPgmodelerUbuntu09openSettings}{0.225}{0.33}%
\locateGraphicTB{14}{width=0.48\paperwidth}{graphics/installingPgModelerUbuntu/installingPgmodelerUbuntu10themeLight}{0.26}{0.25}%
\locateGraphicTB{15}{width=0.48\paperwidth}{graphics/installingPgModelerUbuntu/installingPgmodelerUbuntu11themeLightSet}{0.26}{0.25}%
\locateGraphicTB{16}{width=0.48\paperwidth}{graphics/installingPgModelerUbuntu/installingPgmodelerUbuntu12connectionMode}{0.26}{0.25}%
%
\end{frame}%
%
%
\section{PgModeler unter Microsoft Windows Installieren}%
%
\begin{frame}[t]%
\frametitle{PgModeler unter Microsoft Windows Installieren}%
%
\begin{itemize}%
%
\only<-8>{%
\item Installieren wir nun den \pgmodeler\ unter \microsoftWindows.%
%
\item<2-> Leider ist das nicht direkt möglich.%
%
\item<3-> Es geht, wenn man zuerst das \glsFull{MSYS2} unter \microsoftWindows\ installiert.%
%
\item<4-> \glsShort{MSYS2} ist eine Kollektion von Werkzeugen und Bibliotheken aus der \linux-Welt um eine Umgebung zum Compileren, Installieren und Ausführen von nativer \microsoftWindows-Software zur Verfügung zu stellen\cite{TO2024BCCAIGTMCCAMC:MS}.%
%
\item<5-> Man könnte dann den \pgmodeler\ vom Quelltext compilieren.%
%
\item<6-> Das ist allerdings sehr umständlich.%
%
\item<7-> Wenn wir \glsShort{MSYS2} aber installiert haben, dann können wir den \pgmodeler\ auch über den Paketmanager \bashil{pacman}\cite{VGL2002:P,TOSID2025L:COMPMS} installieren.%
%
\item<8-> Und so werden wir das auch machen.%
}%
%
\only<-10>{%
\item<9-> Wir laden zuerst \glsShort{MSYS2} von der Webseite \url{https://repo.msys2.org/distrib} herunter.%
\item<10-> Wir laden \textil{msys2-x86_64-latest.exe} herunter, wenn wir ein 64~Bit x86~System haben.%
}%
%
\only<-11>{%
\item<11-> Der Download beginnt.%
}%
%
\only<-12>{%
\item<12-> Nachdem der Download fertig ist, führen wir den Installer durch klick auf \menu{Open file} aus.%
}%
%
\only<-13>{%
\item<13-> Im Willkommensbildschrim klicken wir auf \menu{Next}.%
}%
%
\only<-15>{%
\item<14-> Der Installer will \glsShort{MSYS2} in \expandafter\batil{C:\\msys64} installieren.%
\item<15-> Wir stimmen dem zu und klicken auf \menu{Next}.%
}%
%
\only<-16>{%
\item<16-> Wir lassen die Startmenüeinstellungen so wie sie sind und klicken auf \menu{Next}.%
}%
%
\only<-17>{%
\item<17-> Der Installationsprozess beginnt mit dem Entpacken des Archivs.%
}%
%
\only<-19>{%
\item<18-> Wenn der Installer etwa 50\% erreicht, dann kann es passieren, dass er für eine lange Zeit feststeckt.%
\item<19-> Warten Sie dann einfach. Das ist OK. Brechen Sie die Installation \alert{NICHT} ab.%
}%
%
\only<-21>{%
\item<20-> Irgenwann ist die Installation fertig.%
\item<21-> Wir können \inQuotes{Run \pgls{MSYS2} now.} markieren und auf \menu{Finish} klicken.%
}%
%
\only<-22>{%
\item<22-> Von jetzt an können wir \pgls{MSYS2} auch durch Eingabe von \batil{MSYS2} in den Startmenü-Launcher und/oder Klick auf das \pgls{MSYS2}-Icon starten.%
}%
%
\only<-23>{%
\item<23-> Das \glsFull{MSYS2} \glslink{terminal}{Terminal} läuft nun.%
}%
%
\only<-25>{%
\item<24-> \pgls{MSYS2} benutzt den Paketmanager \bashil{pacman}.%
\item<25-> Wir können nach \pgmodeler-Paketen suchen, in dem wir \bashil{pacman -Ss pgmodeler} eintippen und \keys{\return} drücken.%
}%
%
\only<-28>{%
\only<-27>{%
\item<26-> Wir finden sogar mehrere!
}%
\item<27-> Wir wollen das was mit \textil{mingw64} anfängt installieren -- aber nicht das mit plugins.%
\item<28-> Wir selektieren also den Text \bashil{mingw64/mingw64-w64-x86_64-pgmodeler}.%
}%
%
\only<-30>{%
\item<29-> Wir selektieren den Text \bashil{mingw64/mingw64-w64-x86_64-pgmodeler}.%
\item<30-> Dann rechts-klicken wir in das Terminal und klicken auf~\menu{Copy}.%
}%
%
\only<-31>{%
\item<31-> Dann tippen wir \bashil{pacman -S } ein, rechts-klicken in das Terminal und klicken auf~\menu{Paste}.%
}%
%
\only<-33>{%
\item<32-> Insgesamt haben wir also geschrieben \bashil{pacman -S mingw64/mingw64-w64-x86_64-pgmodeler}.%
\item<33-> Wir drücken~\keys{\return}.%
}%
%
\only<-35>{%
\item<34-> Wir sehen eine Liste von Paketen, die installiert und wir werden gefragt, ob wir das wollen.%
\item<35-> Wir tippen auf~\keys{Y} und dann~\keys{\return}.%
}%
%
\only<-37>{%
\item<36-> Die Pakete werden heruntergeladen und installiert.%
\item<37-> Wir kommen zurück zum Terminal prompt.%
}%
%
\only<-39>{%
\item<38-> Wir \alert{schließen} das \glsShort{MSYS2} Terminal und \alert{öffnen} es erneut.%
\item<39-> Wir tippen \bashil{pgmodeler} in das Terminal ein und drücken~\keys{\return}.%
}%
%
\only<-40>{%
\item<40-> Der \pgmodeler\ startet.%
}%
%
\only<-41>{%
\item<41-> Der \pgmodeler\ läuft.%
}%
%
\only<-43>{%
\item<42-> Vielleicht läuft er aber auch nicht.%
\item<43-> Dann müssen wir nochmal in das \pgls{MSYS2} \glslink{terminal}{Terminal} gehen.%
}%
%
\only<-44>{
\item<44-> Wir schreiben \bashil{nano ~/.bashrc} und drücken \keys{\return}.%
}%
%
\only<-46>{
\item<45-> Der Texteditor \bashil{nano} öffnet sich.%
\item<46-> Wir drücken \keys{\ctrl+End}, um bis zum Ende des Textes zu scrollen.%
}%
%
\only<-47>{%
\item<47-> Wir sind am Ende des Textes angekommen.%
}%
%
\only<-48>{%
\item<48-> Wir schreiben \textil{export PATH=\$PATH:/mingw64/bin}.%
}%
%
\only<-50>{%
\item<49-> Wir drücken \keys{\ctrl+O}, um die Datei zu speichern.%
\item<50-> Wenn wir gefragt werden, ob wir den Text speichern wollen, drücken wir \keys{\enter}.%
}%
%
\only<-52>{%
\item<51-> Nachdem wir \keys{\enter} gedrückt haben, wird der Text in die Datei geschrieben.%
\item<52-> Wir drücken nun \keys{\ctrl+X} um den Editor zu verlassen.%
}%
%
\only<-54>{%
\item<53-> Wir sind wieder im \pgls{MSYS2} Terminal.%
}%
\only<-55>{%
\item<54-> Wir schließen das Terminal und öffnen es erneut.%
\item<55-> Danach können wir den \pgmodeler\ ausführen.%
}%
%
\only<-47>{%
\item<56-> Ich verwende \pgmodeler\ im \inQuotes{Light Mode}.%
\item<57-> Um den zu aktivieren, klicken wir auf \menu{\pgmodelerMainMenu>Edit>Settings}.%
}%
%
\only<58>{%
\item<58-> Im Menü \menu{Appearance} ändern wir das \inQuotes{Theme} zu \menu{Light}{\dots}%
}%
%
\only<59>{%
\item Wir klicken auf \menu{Apply}.%
}%
%
\only<60>{%
\item<60-> Im Register \menu{Relationships} stellen wir sicher, dass der \inQuotes{connection mode} auf \inQuotes{Crow's foot notation} gesetzt ist.%
}%
%
\item<61-> Damit haben wir den \pgmodeler\ installiert.%
%
\end{itemize}%
%
\locateGraphicTB{9-10}{width=0.5\paperwidth}{graphics/installingPgModelerWindows/installingPgModelerWindows01msys2Downloads}{0.25}{0.28}%
\locateGraphicTB{11}{width=0.5\paperwidth}{graphics/installingPgModelerWindows/installingPgModelerWindows02msys2DownloadStarts}{0.25}{0.28}%
\locateGraphicTB{12}{width=0.5\paperwidth}{graphics/installingPgModelerWindows/installingPgModelerWindows03msys2Downloaded}{0.25}{0.28}%
\locateGraphicTB{13}{width=0.64\paperwidth}{graphics/installingPgModelerWindows/installingPgModelerWindows04msys2InstallerWelcome}{0.18}{0.29}%
\locateGraphicTB{14-15}{width=0.64\paperwidth}{graphics/installingPgModelerWindows/installingPgModelerWindows05msys2InstallerDest}{0.18}{0.29}%
\locateGraphicTB{16}{width=0.64\paperwidth}{graphics/installingPgModelerWindows/installingPgModelerWindows06msys2InstallerStartMenu}{0.18}{0.29}%
\locateGraphicTB{17}{width=0.64\paperwidth}{graphics/installingPgModelerWindows/installingPgModelerWindows07msys2InstallerUnpacking}{0.18}{0.29}%
\locateGraphicTB{18-19}{width=0.64\paperwidth}{graphics/installingPgModelerWindows/installingPgModelerWindows08msys2InstallerWaiting}{0.18}{0.29}%
\locateGraphicTB{20-21}{width=0.64\paperwidth}{graphics/installingPgModelerWindows/installingPgModelerWindows09runMsys}{0.18}{0.29}%
\locateGraphicTB{22}{width=0.55\paperwidth}{graphics/installingPgModelerWindows/installingPgModelerWindows10manuallyRunMsys}{0.225}{0.29}%
\locateGraphicTB{23}{width=0.64\paperwidth}{graphics/installingPgModelerWindows/installingPgModelerWindows11msysRunning}{0.18}{0.29}%
\locateGraphicTB{24-25}{width=0.64\paperwidth}{graphics/installingPgModelerWindows/installingPgModelerWindows12searchPgmodelerPackages}{0.18}{0.29}%
\locateGraphicTB{26-28}{width=0.55\paperwidth}{graphics/installingPgModelerWindows/installingPgModelerWindows13pgmodelerPackages}{0.225}{0.29}%
\locateGraphicTB{29-30}{width=0.55\paperwidth}{graphics/installingPgModelerWindows/installingPgModelerWindows14copyMingwPgmodeler}{0.225}{0.29}%
\locateGraphicTB{31}{width=0.55\paperwidth}{graphics/installingPgModelerWindows/installingPgModelerWindows15pasteMingwPgmodeler}{0.225}{0.29}%
\locateGraphicTB{32-33}{width=0.55\paperwidth}{graphics/installingPgModelerWindows/installingPgModelerWindows16installCommand}{0.35}{0.3}%
\locateGraphicTB{34}{width=0.55\paperwidth}{graphics/installingPgModelerWindows/installingPgModelerWindows17installYesNo}{0.225}{0.29}%
\locateGraphicTB{35}{width=0.55\paperwidth}{graphics/installingPgModelerWindows/installingPgModelerWindows18installYes}{0.225}{0.29}%
\locateGraphicTB{36}{width=0.55\paperwidth}{graphics/installingPgModelerWindows/installingPgModelerWindows19download}{0.225}{0.29}%
\locateGraphicTB{37}{width=0.55\paperwidth}{graphics/installingPgModelerWindows/installingPgModelerWindows20installed}{0.225}{0.29}%
\locateGraphicTB{38-39}{width=0.64\paperwidth}{graphics/installingPgModelerWindows/installingPgModelerWindows21restartMsys}{0.18}{0.33}%
\locateGraphicTB{40}{width=0.64\paperwidth}{graphics/installingPgModelerWindows/installingPgModelerWindows22pgmodelerLogo}{0.18}{0.27}%
\locateGraphicTB{41}{width=0.64\paperwidth}{graphics/installingPgModelerWindows/installingPgModelerWindows23pgmodelerMain}{0.18}{0.27}%
%
\locateGraphicTB{42-43}{width=0.64\paperwidth}{graphics/installingPgModelerWindows/installingPgModelerWindows24msys2terminal}{0.18}{0.27}%
\locateGraphicTB{44}{width=0.64\paperwidth}{graphics/installingPgModelerWindows/installingPgModelerWindows25nanoBashRc}{0.18}{0.27}%
\locateGraphicTB{45-46}{width=0.64\paperwidth}{graphics/installingPgModelerWindows/installingPgModelerWindows26nanoBashRcOpen}{0.18}{0.27}%
\locateGraphicTB{47}{width=0.64\paperwidth}{graphics/installingPgModelerWindows/installingPgModelerWindows27nanoCtrlEnd}{0.18}{0.27}%
\locateGraphicTB{48}{width=0.64\paperwidth}{graphics/installingPgModelerWindows/installingPgModelerWindows28nanoExportPath}{0.18}{0.27}%
\locateGraphicTB{49-50}{width=0.64\paperwidth}{graphics/installingPgModelerWindows/installingPgModelerWindows29nanoCtrlO}{0.18}{0.27}%
\locateGraphicTB{51-52}{width=0.64\paperwidth}{graphics/installingPgModelerWindows/installingPgModelerWindows30enter}{0.18}{0.27}%
\locateGraphicTB{53-55}{width=0.64\paperwidth}{graphics/installingPgModelerWindows/installingPgModelerWindows31CtrlX}{0.18}{0.27}%
%
\locateGraphicTB{56-57}{width=0.55\paperwidth}{graphics/installingPgModelerUbuntu/installingPgmodelerUbuntu09openSettings}{0.225}{0.33}%
\locateGraphicTB{58}{width=0.48\paperwidth}{graphics/installingPgModelerUbuntu/installingPgmodelerUbuntu10themeLight}{0.26}{0.25}%
\locateGraphicTB{59}{width=0.48\paperwidth}{graphics/installingPgModelerUbuntu/installingPgmodelerUbuntu11themeLightSet}{0.26}{0.25}%
\locateGraphicTB{60-}{width=0.48\paperwidth}{graphics/installingPgModelerUbuntu/installingPgmodelerUbuntu12connectionMode}{0.26}{0.25}%
\end{frame}%
%
\section{PgModeler unter MacOS Installieren}%
%
\begin{frame}[t]%
\frametitle{PgModeler unter MacOS Installieren}%
\begin{itemize}%
%
\item Den \pgmodeler\ unter \pgls{macOS} zu installieren funktioniert so ähnlich wie unter \microsoftWindows.%
%
\item<2-> Weil ich keinen Apple-Computer habe, gibt es diesmal keine Screenshots.%
%
\item<3-> Sie gehen wie folgt vor\only<-3>{.}\uncover<4->{:%
%
\begin{enumerate}%
%
\setcounter{enumi}{0}%
\item Gehen sie zu \url{https://www.macports.org} und dann zum Register \emph{\inQuotes{Installing \pgls{macPorts}}} -- oder direkt zu
\url{https://www.macports.org/install.php}.%
%
\setcounter{enumi}{1}%
\item<5-> Laden Sie die passende \pgls{macPorts}-Version für Ihre Version von \pgls{macOS} herunter\cite{RRJ2008MFUG}.\uncover<6->{ %
Wenn Sie \DEzB\ \inQuotes{\pgls{macOS} Ventura} haben, dann laden Sie den Installer für \inQuotes{\pgls{macOS} Ventura} herunter.}%
%
\setcounter{enumi}{2}%
\item<7-> Installieren Sie der heruntergeladene Version von \pgls{macPorts}.%
%
\setcounter{enumi}{3}%
\item<8-> Updaten Sie die Paketliste des \pgls{macPorts}-Systems\only<-8,14->{.}\only<9-13>{:%
\begin{enumerate}%
%
\item Öffnen Sie ein neues Terminal.%
%
\item<10-> Schreiben Sie \bashil{sudo port -v selfupdate} und drücken Sie~\keys{\enter}.\uncover<11->{ %
Sie werden nach Ihrem (\pgls{sudo}) Passwort gefragt, also Ihrem Passwort für Ihren Computer.\uncover<12->{ %
Geben Sie es ein und drücken~\keys{\return}.%
}}%
%
\item<13-> Nachdem das Update fertig ist, schließen Sie das Terminal.%
\end{enumerate}%
}%
%
\setcounter{enumi}{4}%
\item<14-> Wir brauchen auch \pgls{xcode}.\only<15-17>{ %
Wenn Sie es also noch nicht installiert haben, installieren Sie es jetzt\dots%
\begin{enumerate}%
\item Öffnen Sie ein \emph{neues} Terminal.%
%
\item<16-> Schreiben Sie \bashil{xcode-select --install} und drücken Sie~\keys{\enter}.\uncover<17->{ %
Wenn Sie Ihr Passwort angeben müssen, dann tun Sie das.}%
%
\item<17-> Nachdem die Installation fertig ist, schließen Sie das Terminal.%
\end{enumerate}%
}%
%
\setcounter{enumi}{5}%
\item<18-> Jetzt können wir den \pgmodeler\ installieren\only<-18,23->{.}\only<19-22>{:%
\begin{enumerate}%
\item Öffnen Sie ein neues Terminal.%
%
\item<19-> Schreiben Sie \bashil{sudo port install pgmodeler} und drücken Sie~\keys{\enter}.\uncover<20->{ %
Sie werden nach Ihrem (\pgls{sudo}) Passwort gefragt, also Ihrem Passwort für Ihren Computer.\uncover<21->{ %
Geben Sie es ein und drücken~\keys{\return}.%
}}%
%
\item<22-> Nachdem die Installation fertig ist, schlißen Sie das Terminal.%
\end{enumerate}}%
%
\setcounter{enumi}{6}%
\item<23-> Öffnen Sie ein neues Terminal.%
%
\setcounter{enumi}{7}%
\item<24-> Sie sollten nun den \pgmodeler\ mit dem Kommando \bashil{pgmodeler} + \keys{\enter} ausführen können.%
%
\end{enumerate}}%
\end{itemize}%
\end{frame}%
%
\section{Zusammenfassung}%
%
\begin{frame}%
\frametitle{Zusammenfassung}%
\begin{itemize}%
\item Nun haben wir den \pgmodeler\ installiert.%
\item<2-> Damit können wir logische Modelle für \postgresql-Datenbanken entwerfen.%
\item<3-> Mit dem \pgmodeler\ geht das wesentlich komfortabler als nur mit \glsShort{SQL}.%
\item<4-> Das Coole ist, dass solche Modelle dann direkt nach \glsShort{SQL} übersetzt werden können.%
\item<5-> Und das werden wir bald sehen.%
\end{itemize}%
\end{frame}%
%
\endPresentation%
\end{document}%%
\endinput%
%
